\documentclass[12pt]{article}
\usepackage{ceai}

\begin{document}

\title{CE 488/5588 Homework 07\\
\vspace{2in}
\pmb{Name: \rule{2in}{0.15mm}}
\author{}
}
\maketitle
\newpage

\section*{Topics 11--12: Regression, Clustering, and PCA}

\begin{ch-box}
This homework is based on the lecture content on regression, clustering, and principal
component analysis in Topics 11 and 12. Answer all five questions. Please review the lecture notes or slides while completing the homework.
\end{ch-box}

%\section*{Conceptual Questions and Short Calculations}

\subsection*{Problem 1: Regression versus classification}
Answer the following short questions.

\begin{enumerate}
    \renewcommand{\labelenumi}{(\alph{enumi})}
    \item In one or two sentences, explain the main difference between a classification problem and a regression problem.

    \vspace{18pt}
    \fullunderline
    
    \vspace{18pt} 
    \fullunderline

    \item In classification, a confusion matrix is a standard evaluation tool. Why does this idea usually disappear in regression?

   \vspace{18pt}
    \fullunderline
    
    \vspace{18pt} 
    \fullunderline

    \item Give one reason why regression error is often easier to express mathematically than classification error.

   \vspace{18pt}
    \fullunderline
    
    \vspace{18pt} 
    \fullunderline

    \item Let the regression model be \(\hat{y} = f(\bm{x})\) for samples \((\bm{x}_i,y_i)\), \(i=1,2,\ldots,N\). Write the definitions of the sum of absolute errors (SAE) and the sum of squared errors (SSE).

   \vspace{18pt}
    \fullunderline
    
    \vspace{18pt} 
    \fullunderline
\end{enumerate}

\newpage 

\subsection*{Problem 2: Linear regression with three data points}
Consider the one-dimensional linear regression model
\[
\hat{y} = a_0 + a_1 x.
\]
Use the three data points
\[
(x_1,y_1)=(-1,1), \qquad (x_2,y_2)=(2,2), \qquad (x_3,y_3)=(3,3).
\]

\begin{enumerate}
    \renewcommand{\labelenumi}{(\alph{enumi})}
    \item Write the SAE objective \(E_1(a_0,a_1)\) and the SSE objective \(E_2(a_0,a_1)\) explicitly for these three points.

   \vspace{18pt}
    \fullunderline
    
    \vspace{18pt} 
    \fullunderline

    \item Compute the first-order partial derivatives \(\partial E_2/\partial a_0\) and \(\partial E_2/\partial a_1\).

   \vspace{18pt}
    \fullunderline
    
    \vspace{18pt} 
    \fullunderline

    \item (Gradudate Student only) Form the least-squares system using the conventional augmented design matrix
    \[
    \bm{X}=
    \begin{bmatrix}
    1 & x_1\\
    1 & x_2\\
    1 & x_3
    \end{bmatrix},
    \qquad
    \bm{a}=
    \begin{bmatrix}
    a_0\\
    a_1
    \end{bmatrix},
    \qquad
    \bm{y}=
    \begin{bmatrix}
    y_1\\
    y_2\\
    y_3
    \end{bmatrix}.
    \]
    Substitute the three given points, then compute the least-squares estimate using the pseudo left inverse
    \[
    \bm{X}^{-L}=(\bm{X}^T\bm{X})^{-1}\bm{X}^T,
    \qquad
    \hat{\bm{a}}=\bm{X}^{-L}\bm{y}.
    \]
    Report the numerical values of \(\hat{a}_0\) and \(\hat{a}_1\).

   \vspace{18pt}
    \fullunderline
    
    \vspace{18pt} 
    \fullunderline
\end{enumerate}

\newpage

\subsection*{Problem 3: Practical fitting of a noisy harmonic relation}
Suppose field measurements \((x_i,y_i)\), \(i=1,2,\ldots,N\), are believed to follow the model
\[
y = A\sin(a_0x+b_0) + B\cos(a_1x+b_1) + \varepsilon,
\]
where \(\varepsilon\) represents measurement noise.

Design a practical regression procedure for fitting this model from field data. Your answer should address the following points:

\begin{enumerate}
    \renewcommand{\labelenumi}{(\alph{enumi})}
    \item what data preprocessing or visualization you would perform first;
    \item what objective function you would minimize;
    \item what model (linear, neural network, or support vector regression) would you use to fit the data?
    % \item how you would reduce the risk of overfitting noisy data; and
    \item how you would judge whether the fitted model is satisfactory.
\end{enumerate}

Write your answer as a short, ordered plan rather than as a full essay.

\vspace{24pt}
\fullunderline

\vspace{20pt}
\fullunderline

\vspace{20pt}
\fullunderline

\newpage 

\subsection*{Problem 4: K-means clustering objective and choice of \(K\)}
Let \(\mathcal{D}=\{\bm{x}_1,\bm{x}_2,\ldots,\bm{x}_N\}\) be an unlabeled dataset in \(\mathbb{R}^d\).

\begin{enumerate}
    \renewcommand{\labelenumi}{(\alph{enumi})}
    \item Write the k-means objective function using cluster sets \(S_k\) and centroids \(\bm{\mu}_k\), for \(k=1,2,\ldots,K\).

    \vspace{22pt}
    \fullunderline

    \item State the two alternating steps of the basic k-means algorithm.

    \vspace{20pt}
    \fullunderline

    \item Briefly explain how you would choose \(K\) in each of the following situations:
    \begin{enumerate}
        \renewcommand{\labelenumii}{(\roman{enumii})}
        \item the number of physically meaningful operating regimes is already known from engineering knowledge;
        \item the number of groups is not known in advance.
    \end{enumerate}

    \vspace{20pt}
    \fullunderline
\end{enumerate}

\newpage 


\subsection*{Problem 5: PCA for dimensionality reduction and visualization}
Suppose a dataset is saved in an Excel spreadsheet. Each row corresponds to one instance, and each column corresponds to one measured attribute. Let the data matrix be
\[
\bm{X} \in \mathbb{R}^{N\times d},
\]
where \(N\) is the number of instances and \(d\) is the number of attributes. In this homework, assume \(N=100\) and \(d=1000\).
\begin{enumerate}
    \renewcommand{\labelenumi}{(\alph{enumi})}
    \item Define a centered data matrix \(\bm{X}_c\), and then write the PCA projection from \(\bm{X}_c\) to a reduced three-dimensional representation
    \[
    \bm{Z} \in \mathbb{R}^{100\times 3}.
    \]
    Use a projection matrix \(\bm{W}_3\), and clearly state the dimensions of \(\bm{W}_3\) and \(\bm{Z}\).
    
    \vspace{24pt}
    \fullunderline
    
    \item Briefly explain why PCA can be useful when the spreadsheet contains many attributes, some of which may be noisy, redundant, or not very informative.
    
    \vspace{20pt}
    \fullunderline
    
    \vspace{20pt}
    \fullunderline
    
    \item After computing the first three principal components, imagine a \(3\times 3\) plot grid built from these three coordinates. Describe what should appear in the diagonal panels and what should appear in the off-diagonal panels. Then state one pattern in this plot grid that would suggest meaningful low-dimensional structure in the data.
    
    \vspace{20pt}
    \fullunderline
    
    \vspace{20pt}
    \fullunderline
\end{enumerate}
\end{document}
